\documentclass{article}
\usepackage{physics}
\usepackage{mathrsfs}
\usepackage{cmap}
\usepackage[T1]{fontenc}
\usepackage[utf8]{inputenc}
\usepackage[english,russian]{babel}
\usepackage{indentfirst, setspace}
\usepackage{amsfonts,amssymb}
\usepackage{amsmath,amsthm}
\setlength{\parindent}{0.75cm}

% Кодировки и языки
\usepackage[T2A]{fontenc}
\usepackage[utf8]{inputenc}
\usepackage[english,russian]{babel}

% Математика
\usepackage{amsmath,amssymb,amsfonts}

% Графика и TikZ
\usepackage{graphicx}
\usepackage{tikz}
\usepackage{tikz-3dplot}
\usetikzlibrary{arrows.meta, shapes, positioning, decorations.pathmorphing, patterns, decorations.pathreplacing}

% Оформление оглавления, подписей, шрифта
\usepackage{tocloft}
\usepackage{caption}
\usepackage{tempora}
\usepackage{titlesec}
\usepackage{setspace}
\usepackage{indentfirst}
\usepackage{enumitem}
\usepackage{array}
\usepackage{xcolor}
\usepackage{tabularx}
\usepackage{wrapfig}
\usepackage{hyperref}
\usepackage{listings}
\usepackage{float}

% Настройки списков
\setlist{nosep}
\setlist[itemize]{leftmargin=1.25cm, itemindent=0.65cm}
\setlist[enumerate]{leftmargin=1.25cm, itemindent=0.55cm}

\usepackage[
	a4paper,
	left=30mm,
	right=15mm,
	top=20mm,
	bottom=20mm,
	]{geometry}
	
% Параметры страницы
\geometry{left=30mm,right=10mm,top=20mm,bottom=20mm}
\onehalfspacing
\linespread{1.3}
\parindent=1.25cm

% Нумерация и вид списков
\renewcommand{\theenumi}{\arabic{enumi}}
\renewcommand{\labelenumi}{\arabic{enumi})}
\renewcommand{\theenumii}{.\arabic{enumii}}
\renewcommand{\labelenumii}{\asbuk{enumii})}
\renewcommand{\theenumiii}{.\arabic{enumiii}}
\renewcommand{\labelenumiii}{\arabic{enumi}.\arabic{enumii}.\arabic{enumiii}.}

% Подписи к рисункам и таблицам
\addto\captionsrussian{\renewcommand{\figurename}{Рисунок}}
\DeclareCaptionLabelSeparator{dash}{~---~}
\captionsetup{labelsep=dash}
\captionsetup[figure]{justification=centering,labelsep=dash}
\captionsetup[table]{labelsep=dash,justification=raggedright,singlelinecheck=off}

% Путь к картинкам
\graphicspath{{images/}}

% Маркер для itemize
\def\labelitemi{---}

\frenchspacing

\begin{document}

\selectlanguage{russian}
\fontsize{14}{16}\selectfont

\begin{titlepage}
	\begin{flushleft}
		{\Large\textbf{Министерство образования и науки РФ}}
		
		\vspace{1cm}
		
		{\Large\textbf{Московский государственный технический университет имени Н.Э. Баумана}}
		
		\vspace{1cm}
		
		{\Large\textbf{Факультет <<Фундаментальные науки>>}}
		
		\vspace{1cm}
		
		{\Large\textbf{Кафедра <<Теоретическая механика>> имени профессора\\Н.Е. Жуковского}}
		
		\vspace{3cm}
		
		\textbf{\Large Домашнее задание №\underline{3}}
		
		\vspace{0.4cm}
		
		\underline{Уравнения Лагранжа 2-го рода}\\[-1mm]
		
		{\normalsize \textit{Тема домашнего задания}}
		
		\vspace{1cm}
		
		\begin{tabular}{@{} p{2cm} p{3cm} @{}}
			\textbf{Вариант} & \underline{14} \tabularnewline[4mm]
			\textbf{Группа} & \underline{ФН12-41Б}
		\end{tabular}
		
		\vspace{2cm}
		
		\begin{tabular}{@{} p{3.8cm} p{5.4cm} p{2.5cm} p{3cm} @{}}
			\textbf{Студент} & \underline{Михайлов А.К.} & \hrulefill & \hrulefill \\[-1mm]
			
			& {\normalsize \textit{Фамилия, инициалы}} & {\normalsize \textit{подпись}} & {\normalsize \textit{дата}} \tabularnewline[6mm]
			
			\textbf{Преподаватель} & \underline{Стихно К.А.} & & \\[-1mm]
			
			& {\normalsize \textit{Фамилия, инициалы}} & &
		\end{tabular}
		
		\vfill
		
		\textbf{\textit{Москва, 20\underline{26} год}}
	\end{flushleft}
\end{titlepage}

\hspace{-1.3cm}
\begin{minipage}{0.1\textwidth}
\centering
\begin{tikzpicture}
	
	\coordinate (O) at (0, 0);
	\coordinate (C) at (-4, -1);
	
	\node[above right] at (O) {$O$};
	\node[above left] at (C) {$C$};
	\node[] at (-0.03, 3.05) {$1$};
	\node[] at (2.4, -0.75) {$2$};
	\node[] at (2.3, 1.15) {$3$};
	\node[] at (-5.12, 0.45) {$4$};
	\node[] at (3.4, 2.6) {$\vec{F}$};
	\node[] at (-3.4, 2.45) {$s$};
	\node[] at (-2.45, 3.05) {$A$};
	\node[] at (2.45, 3.05) {$B$};
	\node[] at (-3.7, -2.5) {$\beta$};
	\node[] at (-3.65, 0.7) {$\varphi$};
	\node[] at (-4.1, 1.95) {$O_1$};
	
	\draw[-latex, ultra thick] (-3.8, 2.2) -- (-3, 2.2);
	\draw[-latex, ultra thick] (3, 2.2) -- (3.8, 2.2);
	\draw[thick] (-3.8, 2.17) -- (-3.8, 2.8);
	
	\draw[thick] (0.1, 0) -- (0.25, -0.45);
	\draw[thick] (-0.1, 0) -- (-0.25, -0.45);
	\draw[thick] (-0.4, -0.45) -- (0.4, -0.45);
	\draw[dash dot] (0, 2) -- (0, -2);
	\draw[dash dot] (-2, 0) -- (2, 0);
	\draw[thick] (-0.45, 0.894) -- (-3.96, -0.96);
	\draw[dash dot] (-5, -1) -- (-3, -1);
	\draw[dash dot] (-4, -2) -- (-4, 0);
	
	\draw[thick] (-2.8, 2.5) -- (-2.8, 2.8);
	\draw[thick] (-2, 2.5) -- (-2, 2.8);
	\draw[thick] (-2.8, 2.512) -- (-2, 2.512);
	\fill[pattern=north east lines] (-2.8, 2.512) rectangle (-2, 2.6);
	
	\draw[thick] (2.8, 2.5) -- (2.8, 2.8);
	\draw[thick] (2, 2.5) -- (2, 2.8);
	\draw[thick] (2.8, 2.512) -- (2, 2.512);
	\fill[pattern=north east lines] (2.8, 2.512) rectangle (2, 2.6);
	
	\draw[thick] (2.8, 1.9) -- (2.8, 1.6);
	\draw[thick] (2, 1.9) -- (2, 1.6);
	\draw[thick] (2.8, 1.888) -- (2, 1.888);
	\fill[pattern=north east lines] (2.8, 1.888) rectangle (2, 1.8);
	
	\draw[thick] (-2.8, 1.9) -- (-2.8, 1.6);
	\draw[thick] (-2, 1.9) -- (-2, 1.6);
	\draw[thick] (-2.8, 1.888) -- (-2, 1.888);
	\fill[pattern=north east lines] (-2.8, 1.888) rectangle (-2, 1.8);
	
	\draw[thick] (-3, 2) rectangle (3, 2.4);
	\fill[pattern=north east lines] (-0.4, -0.55) rectangle (0.4, -0.45);
	\fill[pattern=north east lines] (-3.9, 2.17) rectangle (-3.8, 2.8);

	\draw[thick] (O) circle (0.1);
	\draw[thick] (O) circle (1);
	\draw[thick] (O) circle (2);
	\draw[thick] (C) circle (0.05);
	\draw[thick] (C) circle (1);
	\fill (-3.8, 2.2) circle (0.05);
	
	\draw[] (-0.7, 2.3) -- (-0.2, 2.8);
	\draw[] (-0.2, 2.8) -- (0.2, 2.8);
	
	\draw[] (1.7, -0.5) -- (2.2, -1);
	\draw[] (2.2, -1) -- (2.6, -1);
	
	\draw[] (-4.3, -0.2) -- (-4.9, 0.2);
	\draw[] (-4.9, 0.2) -- (-5.3, 0.2);
	
	\draw[] (0.7, 0.25) -- (2.1, 0.9);
	\draw[] (2.1, 0.9) -- (2.5, 0.9);
	
	\draw[] (C) -- (-4, 0.7);
	\draw[] (C) -- (-3.25, 0.5);
	
	\draw[-latex, ultra thick] (-4, 0.4) arc[start angle=90, end angle=63, radius=1.4];
	
	\coordinate (D1) at (-5.3,-2.8);
	\coordinate (D2) at (-2.3,-1.3);
	\coordinate (D3) at (-2.1,-1.3);
	\coordinate (D4) at (-5.1,-2.8);
	\draw[thick, fill=white] (D1) -- (D2) -- (D3) -- (D4) -- cycle;
	\fill[pattern=north east lines] (D1) -- (D2) -- (D3) -- (D4) -- cycle;
	
	\draw[thick] (D1) -- (-1.8, -2.8);
	
	\draw[to-to, thick] (-4,-2.8) arc[start angle=-20, end angle=40, radius=0.5];
\end{tikzpicture}
\end{minipage}
\hfill
\begin{minipage}{0.45\textwidth}
\par
\vspace{5pt}
\textbf{Условие задачи:}
В механической системе рейка 1 массой $m_1$ движется в горизонтальных гладких направляющих под действием постоянной силы $F$. Рейка находится в зацеплении с двухступенчатой шестернёй-барабаном радиуса $R$ барабана 3 радиусом $r$. Радиус инерции шестерни-барабана 3 относительно их оси вращения равен $\rho$, $m_2$ - их общая масса.
\end{minipage}
\par
\vspace{5pt}
\noindent
	К центру C однородного катка 4 массой $m_4$ прикреплена нерастяжимая нить, которая наматывается на барабан 3. Каток 4 катится со скольжением по неподвижной наклонной плоскости с углом наклона $\beta$. К шестерне 3 приложена пара сил с моментом $L_{OZ} = -\alpha \omega_2$ ($\alpha = \text{const} > 0$, $\omega_2$ --- угловая скорость вращения шестерни 2 с барабаном 3). Массой нити и трением качения пренебречь. Принв за обобщенные координаты $q_1 = s$ и $q_2 = \varphi_2$, составить дифференциальные уравнения движения механический системы и помощью уравнений Лагранжа 2-го рода.
\section*{Решение}
\subsection*{1. Анализ механической системы}
Cистема имеет две степени свободы. Обобщённые координаты:
\begin{itemize}
\item $q_1 = s$ --- перемещение рейки 1,
\item $q_2 = \varphi$ --- угол поворота катка 4.
\end{itemize}
\par
\vspace*{0.6cm}
Выразим скорости и угловые скороcти всех частей механизма через обобщенные координаты:
\begin{itemize}
	\item Рейка 1 (масса $m_1$) движется поступательно: 
	\begin{equation*}
		v_1 = \dot{s}
	\end{equation*}
	\item Шестерня 2 и барабан 3 вращаются как единое целое. Так как рейка находится в зацеплении с шестерней без проскальзывания, угловая скорость барабана:
	\begin{equation*}
	\omega_2 = \frac{v_1}{R} = \frac{\dot{s}}{R}
	\end{equation*}
	\newpage
	\item Нить нерастяжима, поэтому поступательная скорость центра масс $C$ катка 4 равна скорости точек на ободе барабана 3:
	\begin{equation*}
	v_C = \omega_2 r = \frac{r}{R} \dot{s}
	\end{equation*}
	\item Каток 4 катится со скольжением по неподвижной наклонной плоскости. Его угловая скорость является независимой координатой:
	\begin{equation*}
	\omega_4 = \dot{\varphi}
	\end{equation*}
\end{itemize}
\subsection*{2. Вычисление кинетической энергии системы}
Кинетическая энергия системы $T$ равна сумме кинетических энергий её частей: $T = T_1 + T_{23} + T_4$.
\begin{itemize}
	\item Кинетическая энергия поступательно движущейся рейки:
	$$ T_1 = \frac{m_1 v_1^2}{2} = \frac{m_1 \dot{s}^2}{2} $$
	\item Кинетическая энергия вращающихся шестерни 2 и барабана 3. Момент инерции этой подсистемы относительно оси вращения равен $I_{Oz} = m_2 \rho^2$ (где $\rho$ --- радиус инерции).
	$$ T_{23} = \frac{I_{Oz} \omega_2^2}{2} = \frac{m_2 \rho^2}{2} \left( \frac{\dot{s}}{R} \right)^2 = \frac{m_2 \rho^2}{2 R^2} \dot{s}^2 $$
	\item Каток 4 совершает плоское движение. Примем его за однородный цилиндр радиуса $r_4$, тогда его момент инерции относительно центра масс $I_C = \frac{m_4 r_4^2}{2}$. По теореме Кёнига:
	$$ T_4 = \frac{m_4 v_C^2}{2} + \frac{I_C \omega_4^2}{2} = \frac{m_4}{2} \left( \frac{r}{R} \dot{s} \right)^2 + \frac{m_4 r_4^2}{4} \dot{\varphi}^2 = \frac{m_4 r^2}{2 R^2} \dot{s}^2 + \frac{m_4 r_4^2}{4} \dot{\varphi}^2 $$
\end{itemize}
\par
Суммируя эти компоненты, получаем полную кинетическую энергию:
\begin{equation*}
T = \frac{1}{2} \left( m_1 + \frac{m_2 \rho^2 + m_4 r^2}{R^2} \right) \dot{s}^2 + \frac{m_4 r_4^2}{4} \dot{\varphi}^2
\end{equation*}
\par
Для подстановки в уравнения Лагранжа найдём необходимые частные и полные производные:
\begin{equation*}
\frac{\partial T}{\partial s} = 0, \quad \frac{\partial T}{\partial \varphi} = 0
\end{equation*}
\newpage
\begin{equation*}
\frac{\partial T}{\partial \dot{s}} = \left( m_1 + \frac{m_2 \rho^2 + m_4 r^2}{R^2} \right) \dot{s} \implies \frac{d}{dt} \left( \frac{\partial T}{\partial \dot{s}} \right) = \left( m_1 + \frac{m_2 \rho^2 + m_4 r^2}{R^2} \right) \ddot{s}
\end{equation*}
\begin{equation*}
\frac{\partial T}{\partial \dot{\varphi}} = \frac{m_4 r_4^2}{2} \dot{\varphi} \implies \frac{d}{dt} \left( \frac{\partial T}{\partial \dot{\varphi}} \right) = \frac{m_4 r_4^2}{2} \ddot{\varphi}
\end{equation*}
\subsection*{3. Определение обобщённых сил}
\hspace{2.3cm}
\begin{tikzpicture}[]
	
\coordinate (O) at (0, 0);
\coordinate (C) at (-4, -1);

\node[above right] at (O) {$O$};
\node[font=\small, left] at (C) {$C$};
\node[] at (3.4, 2.6) {$\vec{F}$};
\node[] at (-3.75, 2.5) {$s$};
\node[] at (-3.24, 2.55) {$\delta s$};
\node[] at (-2.45, 3.05) {$A$};
\node[] at (2.45, 3.05) {$B$};
\node[] at (-3.7, -2.5) {$\beta$};
\node[] at (-3.65, 0.7) {$\varphi$};
\node[] at (-2.9, 0.45) {$\delta\varphi$};
\node[font=\small] at (1.5, 0.7) {$L_{OZ}$};
\node[] at (-5.1, -2.3) {$F_{\text{тр}}$};
\node[] at (-4.6, 1.95) {$O_1$};
\node[font=\small] at (-4.45, -1.55) {$m_4\vec{g}$};
\node[font=\small] at (-0.41, -0.63) {$m_2\vec{g}$};
\node[] at (-0.7, -2.45) {$\omega_2$};
\node[] at (-5.4, -0.2) {$\omega_4$};
\node[] at (-4.6, -0.6) {$\vec{N}$};

\draw[-latex, ultra thick] (-4.4, 2.2) -- (-3.5, 2.2);
\draw[-latex, ultra thick] (-3.5, 2.2) -- (-3, 2.2);
\draw[-latex, ultra thick] (3, 2.2) -- (3.8, 2.2);
\draw[thick] (-4.4, 2.17) -- (-4.4, 2.8);
\draw[thick] (-3.5, 2) -- (-3.5, 2.4);

\draw[thick] (0.1, 0) -- (0.25, -0.45);
\draw[thick] (-0.1, 0) -- (-0.25, -0.45);
\draw[thick] (-0.4, -0.45) -- (0.4, -0.45);
\draw[thick] (-0.45, 0.894) -- (-3.96, -0.96);

\draw[thick] (-2.8, 2.5) -- (-2.8, 2.8);
\draw[thick] (-2, 2.5) -- (-2, 2.8);
\draw[thick] (-2.8, 2.512) -- (-2, 2.512);
\fill[pattern=north east lines] (-2.8, 2.512) rectangle (-2, 2.6);

\draw[thick] (2.8, 2.5) -- (2.8, 2.8);
\draw[thick] (2, 2.5) -- (2, 2.8);
\draw[thick] (2.8, 2.512) -- (2, 2.512);
\fill[pattern=north east lines] (2.8, 2.512) rectangle (2, 2.6);

\draw[thick] (2.8, 1.9) -- (2.8, 1.6);
\draw[thick] (2, 1.9) -- (2, 1.6);
\draw[thick] (2.8, 1.888) -- (2, 1.888);
\fill[pattern=north east lines] (2.8, 1.888) rectangle (2, 1.8);

\draw[thick] (-2.8, 1.9) -- (-2.8, 1.6);
\draw[thick] (-2, 1.9) -- (-2, 1.6);
\draw[thick] (-2.8, 1.888) -- (-2, 1.888);
\fill[pattern=north east lines] (-2.8, 1.888) rectangle (-2, 1.8);

\draw[thick] (-3, 2) rectangle (3, 2.4);
\fill[pattern=north east lines] (-0.4, -0.55) rectangle (0.4, -0.45);
\fill[pattern=north east lines] (-4.5, 2.17) rectangle (-4.4, 2.8);

\draw[thick] (O) circle (0.1);
\draw[thick] (O) circle (1);
\draw[thick] (O) circle (2);
\draw[thick] (C) circle (0.05);
\draw[thick] (C) circle (1);
\fill (-4.4, 2.2) circle (0.05);

\draw[] (C) -- (-4, 0.7);
\draw[] (C) -- (-3.25, 0.5);
\draw[] (C) -- (-2.8, 0.2);

\draw[-latex, ultra thick] (-4, 0.4) arc[start angle=90, end angle=63, radius=1.4];
\draw[-latex, ultra thick] (-3.38, 0.26) arc[start angle=63, end angle=44, radius=1.4];

\coordinate (D1) at (-5.3,-2.8);
\coordinate (D2) at (-2.3,-1.3);
\coordinate (D3) at (-2.1,-1.3);
\coordinate (D4) at (-5.1,-2.8);
\draw[thick, fill=white] (D1) -- (D2) -- (D3) -- (D4) -- cycle;
\fill[pattern=north east lines] (D1) -- (D2) -- (D3) -- (D4) -- cycle;

\draw[thick] (D1) -- (-1.8, -2.8);

\draw[to-to, thick] (-4,-2.8) arc[start angle=-20, end angle=40, radius=0.5];

\draw[-latex, ultra thick] (O) -- (0, -0.9);
\draw[-latex, ultra thick] (C) -- (-4, -1.9);
\draw[-latex, ultra thick] (-3.55, -1.9) -- (-4.4, -0.2);
\draw[-latex, ultra thick] (-3.55, -1.9) -- (-4.7, -2.49);

\draw[-latex, ultra thick] (1.2, 0) arc[start angle=0, end angle=50, radius=1.2];
\draw[-latex, ultra thick] (0, -2.2) arc[start angle=-90, end angle=-120, radius=2.2];
\draw[-latex, ultra thick] (-5.2, -1) arc[start angle=-180, end angle=-223, radius=1.2];

\end{tikzpicture}

Вычислим обобщённые силы $Q_s$ и $Q_\varphi$, соответствующие координатам $s$ и $\varphi$. Внешние и внутренние связи, наложенные на систему (реакции опор и направляющих), работы не совершают, так как являются идеальными. Предполагаем, что скорость проскальзывания точки контакта катка направлена вверх по наклонной плоскости, следовательно, вектор силы трения направлен вниз по плоскости.
\par
\vspace{0.2cm}
а) Пусть $\delta s \neq 0$ и $\delta \varphi = 0$. При этом работу совершают постоянная сила $\vec{F}$, момент $L_{OZ}$, сила тяжести катка и сила трения скольжения.

Сумма элементарных работ всех активных сил на этом перемещении:
$$ 
\delta A_s = F \delta s + L_{Oz} \delta \varphi_2 - m_4 g \sin\beta \delta s_C - F_{\text{тр}} \delta s_C
$$
\par
Подставим заданный момент сопротивления $L_{Oz} = -\alpha \omega_2 = -\alpha \frac{\dot{s}}{R}$:
$$
\delta A_s = F \delta s - \alpha \frac{\dot{s}}{R^2} \delta s - m_4 g \sin\beta \frac{r}{R} \delta s - F_{\text{тр}} \frac{r}{R} \delta s
$$
\par
Отсюда обобщённая сила, соответствующая координате $s$:
$$
Q_s = \frac{\delta A_s}{\delta s} = F - \frac{\alpha \dot{s}}{R^2} - \frac{r}{R} m_4 g \sin\beta - \frac{r}{R} F_{\text{тр}}
$$
\par
\vspace{0.2cm}
б) Теперь пусть $\delta \varphi \neq 0$ и $\delta s = 0$ (центр масс катка неподвижен). На этом перемещении совершает работу только сила трения скольжения, которая создаёт момент относительно центра масс катка:
$$ 
\delta A_\varphi = F_{\text{тр}} r_4 \delta \varphi \implies Q_\varphi = \frac{\delta A_\varphi}{\delta \varphi} = F_{\text{тр}} r_4
$$
\par
\vspace{0.2cm}
Осталось определить силу трения $F_{\text{тр}}$. По условию каток 4 скользит по наклонной плоскости, значит, сила трения $F_{\text{тр}}$ является силой трения скольжения, величина которой пропорциональна нормальной реакции опоры $N$: 
$$
F_{\text{тр}} = \mu N
$$
\par
Чтобы найти $N$, запишем теорему о движении центра масс катка 4 в проекции на ось, перпендикулярную наклонной плоскости:
$$ 
N - m_4 g \cos\beta = 0 \implies N = m_4 g \cos\beta. 
$$
\par
Следовательно, сила трения скольжения:
$$
F_{\text{тр}} = \mu m_4 g \cos\beta. 
$$

Подставим полученную силу трения в выражения для обобщённых сил:
$$ Q_s = F - \frac{\alpha}{R^2} \dot{s} - \frac{r}{R} m_4 g (\sin\beta + \mu \cos\beta), $$
$$ Q_\varphi = \mu m_4 g r_4 \cos\beta. $$

\subsection*{4. Составление уравнений Лагранжа 2-го рода}
Дифференциальные уравнения движения механической системы в форме уравнений Лагранжа 2-го рода имеют вид:
$$ 
\frac{d}{dt} \left( \frac{\partial T}{\partial \dot{q}_i} \right) - \frac{\partial T}{\partial q_i} = Q_i
$$
\par
Подставляя найденные производные кинетической энергии и обобщённые силы, получаем итоговую систему дифференциальных уравнений:
\begin{equation*}
	\begin{cases}
		\left( m_1 + \frac{m_2 \rho^2 + m_4 r^2}{R^2} \right) \ddot{s} = F - \frac{\alpha}{R^2} \dot{s} - \frac{r}{R} m_4 g (\sin \beta + \mu \cos \beta) \\
		\ddot{\varphi} = \frac{2 \mu g \cos \beta}{r_4}
	\end{cases}
\end{equation*}

\end{document}
